\resetdatestamp

\newcommand\Dfrac[2]{\frac{\displaystyle #1}{\displaystyle #2}}
\newcommand{\mathBF}[1]{\mbox{\boldmath $#1$}}
\newcommand{\C}[1]{\mathBF{#1}}

\chapter{Mapping}

\begin{comment}
\TeX{} does a marvelous job of setting mathematical formulas, most often
 choosing pleasing spacing.
However, on occasion one should intercede to improve the layout.
This chapter defines a few such occasions.
In addition, this chapter documents some features of the {\tt amsmath}
 package which overcome difficulties in typesetting some mathematical
 forms.
The {\tt amsmath} package is documented 
 in {\it The \LaTeX{} Companion} \cite{Goossens:1997}.

The modified setup is typeset as
\begin{equation}
  G(z) = \begin{cases}
           \Dfrac {P(z)}{1+z^{-1}} & \text{for $p$ even}, \\[1ex]
           P(z)                    & \text{for $p$ odd}.
         \end{cases}
\end{equation}

With the modified definitions, we get the following.
\def\hC#1{\C{\hat{#1}}\vphantom{\C{#1}}}           % hat vector
\def\htC#1{\C{\hat{\tilde{#1}}}\vphantom{\C{#1}}}  % hat, tilde vector
\def\tC#1{\C{\tilde{#1}}\vphantom{\C{#1}}}         % tilde vector
\begin{equation}
\begin{split}
  \C{d}^{(i)} &= \hC{v}^{(i)} - \htC{v}^{(i)} \\
  \C{n}^{(i)} &= \C{u}^{(i)} - \tC{v}^{(i)}
\end{split}
\end{equation}
\end{comment}

\section{Definition}

\begin{comment}
texture mapping
conceptual metaphor
function
morphism
\end{comment}

The term \emph{mapping} is used in mathematics, computer science, and related technical fields to encapsulate the concept of a \emph{function} and other related concepts under one abstraction. Like most technical terms used in multiple technical disciplines, the meaning of \emph{mapping} varies according to the context of the field that it is used. However in the case of mapping, the variation is subtle enough that a cursory understanding of the way mapping is used in a couple of contexts outside of DMIs might illuminate why it is a necessary and valuable concept to encapsulate in the context of new musical interfaces and help understand what is required for an interface designed specifically for mapping.

\subsection{Mappings in Mathematics}
\label{sec:Mappings in Mathematics}

In the mathematics context, the use of the term \emph{mapping} varies subtly according to the sub-discipline within mathematics in which it is used and often between specific mathematicians; however it is acceptable to use the word as a synonym for \emph{function}. A function in mathematics describes an associative relationship between two sets of numbers. One set of numbers is referred to as a \emph{domain} or \emph{set of inputs} and the other set is referred to as a \emph{codomain}, \emph{range}, or \emph{set of outputs}. Mathematicians often say that a function "maps" a domain to or onto a range, hence the reason for using the term \emph{mapping} as a synonym for \emph{function} \cite{functionMapping}.

Strictly speaking in mathematics, a function can only represent a \emph{one-to-one mapping} or a \emph{many-to-one mapping}. For now it is helpful to know that a property of both types of mapping is that every element in the input (domain) set is associated by the function with one and only one element in the output (range) set. A corollary statement about a one-to-one and many-to-one mappings is that every element in the output set might be associated to one element (one-to-one) or many elements (many-to-one) in the input set. This language should become more clear as the usage of the terms \emph{one-to-one} and \emph{many-to-one} are explored in the remainder of the chapter.

The following are examples of equations that represent functions:

\begin{equation} \label{eq1}
y = f(x) = x + 1
\end{equation}
\begin{equation} \label{eq2}
z = g(x) = x^{2}
\end{equation}

In addition to both being functions, equation \ref{eq1} is a one-to-one mapping and equation \ref{eq2} is a many-to-one mapping. Basic types of functions like equations \ref{eq1} and \ref{eq2} often use the set of all real numbers (roughly, defined as any number that can be expressed as a potentially infinite string of digits - in the case of irrational numbers like \begin{math}\pi\end{math} - in decimal notation) as both the input set and the output set. Since the set of all real numbers is infinitely large, the association from one input element to one output element is described by a mathematical transformation that when applied to any element of the input set (the set of all real numbers) produces the element of the output set (the set of all real numbers) that the input element is mapped to. This is convenient because that means the function can be written very concisely in one line as opposed to two infinitely long columns with lines drawn between elements that are mapped to each other (not very convenient). 

Equation \ref{eq1} is the canonical notation for representing a mapping/function labeled \begin{math}f\end{math} where the input element is labeled \begin{math}x\end{math} and the output element is labeled \begin{math}y\end{math}. Equation \ref{eq2} represents a mapping/function labeled \begin{math}g\end{math} where the input element is labeled \begin{math}x\end{math} and the output element is labeled \begin{math}z\end{math}. 

To calculate the mapping that is implied by these two mathematical transformations, substitute any real number (since the set of all real numbers is the input set) for the symbol \begin{math}x\end{math}. By this logic, equation \ref{eq1} represents the mapping where every real number \begin{math}x\end{math} is mapped to the real number \begin{math}y\end{math} that is the value of \begin{math}x + 1\end{math}. Equation \ref{eq2} represents the mapping where every real number \begin{math}x\end{math} is mapped to the real number \begin{math}z\end{math} that is the value of \begin{math}x^{2}\end{math}. Equation \ref{eq2} is a many-to-one mapping because it maps both the positive and negative of every input number to the square of that number; unlike equation \ref{eq1} which maps every input number to an output number that no other input number is mapped to.

\subsection{Cardinality in Mathematics}

There is more to be said about the concept of a one-to-one mapping and two related concepts, one-to-many mappings and many-to-one mappings. The framework for discussing different types of mappings is provided in one form by the mathematical concept of \emph{cardinality}.

First, three additional concepts should be enumerated and elaborated upon. These concepts are \emph{injection}, \emph{surjection}, and \emph{bijection}.

An \emph{injection} is a function that meets the condition that every element of the function's codomain (output) set is associated to no more than one element in the functions domain (input) set. Here, domain and codomain are used in the same way as in section \ref{sec:Mappings in Mathematics}.

A \emph{surjection} is a function that meets the condition that every element in the codomain set is associated to at least one element in the domain set.

A \emph{bijection} is a function that is both an \emph{injection} and a \emph{surjection}.

\begin{figure}[ht]
\centering
\subfloat[Subfigure 1 list of figures text][Non-injective and non-surjective]{
\includegraphics[width=0.35\textwidth]{200px-Total_function.png}
\label{fig:figure1}}
\qquad
\subfloat[Subfigure 2 list of figures text][Injection]{
\includegraphics[width=0.35\textwidth]{200px-Injection.png}
\label{fig:figure2}}
\qquad
\subfloat[Subfigure 3 list of figures text][Surjection]{
\includegraphics[width=0.35\textwidth]{200px-Surjection.png}
\label{fig:figure3}}
\qquad
\subfloat[Subfigure 4 list of figures text][Bijection]{
\includegraphics[width=0.35\textwidth]{200px-Bijection.png}
\label{fig:figure4}}
\caption{Four types of functions}
\label{fig:globfig}
\end{figure}

Figure \ref{fig:globfig} summarizes the possible types of functions using these distinctions.

Within mathematics, cardinality is a concept that allows one to reason about the relative sizes of different sets. It is particularly useful when comparing infinite sets. It allows one to reason correctly, for example, that the infinite size of the set of all real numbers is larger than the set of all positive integers (natural numbers).

Two sets are defined to have the same \emph{cardinality} if there exists a bijection between the two sets. It follows from the definition of injection that if there exists an injection from set \emph{A} to set \emph{B}, but no bijection, then the cardinality of B is greater than the cardinality of A \cite{cardmath2010}.

As mentioned before, these mathematical concepts are useful for describing one-to-one mappings and many-to-one mappings. In this terminology, injections and bijections are one-to-one mappings, whereas surjections and functions that are neither injections or surjections are many-to-one mappings.

\subsection{Mappings in Computer Science}
\label{sec:Mappings in Computer Science}

In the computer science context, a mapping is an abstract data type or concrete data structure, commonly referred to as an \emph{associative array} or \emph{dictionary} \cite{assocarray2008}.  

Abstractly and simply, the one and only way to store data of any kind (the value "0", the value "1", some text, an image, a sound, a recorded signal) in a computer of any kind, such that one can reliably retrieve the data later, is to put the data in some "bucket" that has some "label". The words "bucket" and "label", as used here, are not well-defined in computer science. But they help to unobtrusively make the point that if the data has no label that is associated with it or equivalently if the label has no unique memory location associated with it, then practically, it cannot be retrieved. In the scenario of a computer processor accessing internal memory, \emph{bucket} and \emph{label} could be replaced by \emph{memory location} and \emph{memory address} respectively.

As an example from the music information retrieval context, if one wants to be able to store a large amount of sheet music such that it can be found later, one needs some system to catalog or index the material such that a piece of sheet music can be referred to by some label. In the case of a library, this label is often a call number, the title of the piece, the name of the composer, etc.

Just like a reference library can use many different labeling systems to catalog the sheet music stored within its walls so that it can be retrieved with minimum headaches, a computer system can internally utilize many different abstract data types or concrete data structures. In both cases, the best choice is dependent on the material or data being cataloged or indexed. A mapping as a data structure in computer science parlance then, can be understood as a particular type of system for labeling data.

In particular, a mapping is a labeling system that is composed of a set of labels (keys) and a set of pieces of data (values), where every label in the first set is associated with one or many of the pieces of data in the second set. Then this structure encompassing the two sets and the list of associations between the two sets is given a higher level label that refers to the mapping as a whole. This allows one to construct a mapping of mappings, which lends itself nicely to hierarchical labeling systems.

This reframing of the term in the computer science context, as opposed to the mathematical context, focuses on one implication of the concept of mapping. One can specify multiple, unique mappings that each operate on the same collection of labels and the same collection of pieces of data. Similarly, one can specify multiple, unique functions that each operate on the same input set and the same output set. 

So in a computer system, one can create multiple mappings that each associate the same set of labels with the same set of data differently. This characteristic of mappings is one of the primary reasons why mappings are a useful frame for understanding the potential of digital musical instruments. 

\subsection{Cardinality in Computer Science}

Within computer science, cardinality is a concept used in relational database systems \cite{carddata1987}. 

A relational database can be seen as a very robust and flexible \emph{data labeling system}, in reference to the discussion in section \ref{sec:Mappings in Computer Science}. A relational database is defined as a set of relations, where a relation is defined as a set of lists of attributes. Each list in the set has the same attributes as every other list in the set. An attribute defines a property that can take any value constrained by the set of values that are defined for each particular attribute. A relation is usually described by a table as in Table \ref{tab:studentrelation}.

\begin{table}
    \begin{center}
    \begin{tabular}{ | l | l | l | l | }
    \hline
    Student ID : Integer & Name : String & Thesis Topic : String \\ \hline
    00001 & Alice & 4d Maps \\ \hline
    00002 & Bob &  Anti-Gravity Toaster \\ \hline
    00003 & Carl & Invisible Clothing \\ \hline
    00004 & Doug & Personal Babel Fish \\
    \hline
    \end{tabular}
    \end{center}
    \caption{Example of a database relation - Student Table}
    \label{tab:studentrelation}
\end{table}

\begin{table}
    \begin{center}
    \begin{tabular}{ | l | l | }
    \hline
    Faculty ID : Integer & Name : String \\ \hline
    0001 & Ohyes Hedid \\ \hline
    0002 & Ohno Hedidnt \\
    \hline
    \end{tabular}
    \end{center}
    \caption{Example of a database relation - Faculty Table}
    \label{tab:facultyrelation}
\end{table}

\begin{table}
    \begin{center}
    \begin{tabular}{ | l | l | l | }
    \hline
    Pairing ID : Integer & Faculty ID : Integer & Student ID : Integer \\ \hline
    00000001 & 0001 & 00004 \\ \hline
    00000002 & 0002 & 00002 \\ \hline
    00000003 & 0001 & 00003 \\ \hline
    00000004 & 0001 & 00001 \\ 
    \hline
    \end{tabular}
    \end{center}
    \caption{Example of a database relation - Pairing Table}
    \label{tab:relation}
\end{table}

In a relational database, the cardinality of an attribute is the number of instances of that attribute that can be associated with a given set of instances for other attributes. In a database as in DMI mapping, it is largely sufficient to distinguish between a cardinality of \emph{one} and a cardinality of \emph{many}.

The simple database composed of the relations in Tables \ref{tab:studentrelation}, \ref{tab:facultyrelation}, and \ref{tab:relation} provides an example to reason about the cardinality of database relations. If every student (represented by their ID) has one advisor (also represented by their ID) but every advisor has many students, then the relation is one-to-many from Advisor ID to Student ID. Since the horizontal ordering of attributes has no meaning in a relation, we can also reason that the relation from Student ID to Advisor ID is the reverse, many-to-one. Finally, every student has one thesis topic and each thesis topic is the topic of one student so the relation is one-to-one from Student ID to Thesis Topic.

These mathematical and computer science concepts provide a solid foundation for understanding the various nuances of mapping abstractly and concretely. This background will be necessary to create a well informed user interface for creating and modifying DMI mappings.

\section{Mapping in Digital Musical Instruments}
\label{sec:Importance of Mapping}

In an acoustic musical instrument, the equivalents of the "control mechanism" and "sound generation mechanism", as they are called in a DMI, are \emph{intrinsically coupled} or bound to each other because the physical material (e.g. wood, metal, horse hair) of the instrument that forms the control mechanism also forms the sound generation mechanism. Making a decision that alters the control mechanism will significantly or subtly alter the sound generation mechanism. Replacing the horse hair of a violin bow with artificial fiber will result in a perceived difference in tone to the trained ear.

In a DMI, the control mechanism is composed of sensors typically embedded in structural materials that form the shape of the instrument and connected to a microcontroller that provides the interface for accessing the sensor signals. The physical gestures of the performer are not coupled to the sound generation through purely mechanical means, but through electrical means. Each sensor, whether it senses acceleration, orientation, touch, etc. produces an electrical signal that correlates to some physical measurement. This analog electrical signal is converted to a digital signal through some kind of analog to digital signal conversion (perhaps within the control mechanism on an integrated microcontroller) and then this digital signal is processed by a computer and sent to the sound generation mechanism, which is a piece of software that produces audio samples.

The following illustrative example makes clear that the control mechanism and sound generation mechanism are not \emph{intrinsically coupled} in a DMI as they are in an acoustic musical instrument.

If one takes a control mechanism with embedded touch sensors, records an output signal produced by a performer with the control mechanism on a hypothetical signal recorder, and replaces the connection between the control mechanism and the sound generation mechanism with a connection between the signal recorder and the sound generation mechanism, it is clear that the active control mechanism of the DMI has changed. The control mechanism is now the signal recorder. The original control mechanism with the touch sensors is no longer part of the DMI system because it is disconnected from the DMI. But this change results in no change within the sound generation mechanism. The sound is generated by software that produces audio samples and thus is not affected by a change to the physical hardware of the control mechanism as long as the control mechanism produces an output signal that is identical. If the performer presses play on the recorder it will be impossible to tell the difference between the initial DMI and the modified DMI purely from the sound generated by the DMI.

In physical terms, when two formerly intrinsically coupled aspects of system are free to become uncoupled, a new degree of freedom is introduced. One imprecise way to understand why this happens is that the "signal" (used loosely) controlling the sound generation is no longer all the vibrations in the physical body of the instrument (in an acoustic musical instrument), but a discrete electrical signal that travels through a very specific electrical channel (in a DMI).

In the context of DMI research this new degree of freedom is the freedom to choose a mapping \cite{hunt2003}. Mapping as used in this context is very similar to how the term is used in both the mathematics and computer science context. Each context provides an understanding that is well suited to manipulate one of two possible conceptual layers in a mapping system. These two contexts have also been referred to as a \emph{systems} point of view and a \emph{functional} point of view respectively \cite{nort2010}. 

\begin{enumerate}
\item The set of sensors embedded in the control mechanism outputs a set of signals, which we will refer to as \emph{output signals}. The sound generation mechanism outputs sound that is controlled through a set of \emph{input signals} defined by the software that is responsible for the sound generation. Therefore, for any given control mechanism and sound generation mechanism there is a choice of how to associate the two sets of signals (output signals from the control mechanism and input signals from the sound generation mechanism) with each other much like an associative array in the computer science context. Importantly, this is a choice that is not independently available when constructing, for example, a violin because the two sets of signals (if this distinction between control mechanism and sound generation mechanism is imagined in a violin) in a violin are physically coupled and inseparable. It is not a choice that is explicitly made when constructing an acoustic instrument.

\item Furthermore, there is a second lower level form of mapping when dealing with a specific pair of associated signals. Because it is unlikely that any given signal output by a control device is calibrated to a particular input signal in a sound generation device that one might associate with the given output signal, we require a second set of relations beyond a binary associated/not associated. Any given output signal might be capable of outputing values that are outside the set of values that a particular input handles or simply be outputting a different data type altogether. In electronics, this is called \emph{signal conditioning} and can be handled completely in the analog components that produce the signals. If this signal conditioning is not handled in the analog hardware and is modified dynamically in the software of the DMI after analog to digital conversion of the signal then this signal conditioning is handled by the mapping system. In this case, a mapping in the functional, mathematical sense of the word is necessary to transform the output signal so that the values reliable fall within the set of acceptable input values. Performing signal conditioning as part of the DMI mapping, as opposed to as part of the control mechanism, has the benefit of not requiring any change in the analog electronics of the DMI if the associative mapping layer
\end{enumerate}

A prototypical mapping system can specify not only the associations between two sets of signals (pairs of connected signals), but associate a function with each pair of connected signals that specifies how to transform/condition the signal before copying the value to an input. This function might in fact be simply \begin{math}y = x\end{math} in the case where no calibration is needed, however the point is it is not practical to assume that this will always be an acceptable mapping function between an output and input and signal conditioning might be included as a second layer in the a mapping system.

This physical decoupling between the control mechanism, the mapping system, and the sound generation mechanism creates a modularity between the three components that can be exploited to simplify the DMI development and maintenance process.

\section{A Dedicated Mapping System: Libmapper}

The creation, evolution, and stabilization of the design of a DMI is inherently a tightly iterative process. A violin luthier must be comfortable enough playing a violin to develop an aesthetic sense of what effect the countless decisions that are made in the process of building a violin have on the experiential quality of the instrument that they are creating. Whether a particular decision is the right one can only truly be evaluated after the decision is made and someone plays the violin. 

Similarly, a team of people building a DMI will often make a construction decision, evaluate the effect of that decision, make some changes, and repeat. This is often referred to as iterative development and occurs in any serious design process whether or not it is ever made explicit \cite{iterative2003}. 

Because such iterative development is so beneficial to the development of DMIs, people often use various platforms and frameworks as common foundational building blocks for such systems. These environments often take care of implementing higher level functionality for controlling and generating audio in realtime or interfacing with hardware and thus allow one to quickly test ideas for a DMI and focus on experimenting with the aspects of the control mechanism and sound generation that are unique to the particular project before investing too much time and money into custom parts without assurance that the ideas are viable in a basic sense. For the control hardware there are platforms like Arduino. Coding environments like Max/MSP and code libraries like STK simplify the development of synthesizer modules for sound generation. In an effort to provide a similar foundation for creating mapping systems, the Digital Orchestra Toolbox was created as part of the McGill Digital Orchestra project and includes components representing the common subroutines used in a mapping system. More recently, this functionality has been reimplemented as a C library called Libmapper at the IDMIL at McGill University. Libmapper has several characteristics which make it a desirable system for configuring mappings. 

One is that a mapping between control and sound generation can be modified without recompiling any code, restarting any system, or reloading any script. Typically, if a mapping is specified in software through a Max/MSP program, C program, etc. the part of the code that specifies the mapping must be modified and recompiled or reinterpreted before the new mapping becomes active. Any DMI that embeds Libmapper in its control and sound generation software can have its mapping modified simply by being sent specific messages as specified by the mapping protocol. 

Another useful characteristic is that in the case of a local network with many different control mechanisms and sound generation mechanisms available, no central device is needed to facilitate communication between devices. If one component on the network expriences some difficulties, the other devices will still be able to keep the mappings between the remaining devices operating.  

Lastly, and of particular relevance to the primary topic of this paper, is that the current state of mappings between various devices on the network can be viewed and modified by any graphical user interface that embeds Libmapper. In Libmapper parlance, these applications can register as \emph{monitors}. Because of this capability, multiple members of a team can be viewing and modifying the same mappings for multiple DMIs using different computers and using different graphical user interfaces.  

A network with Libmapper-enabled devices and monitors is a \emph{Mapper network}.

\subsection{Open Sound Control Protocol}

Seen from a bird's eye conceptual level, Libmapper is an implementation of a communication protocol that was created for specifying mappings over a distributed network \cite{Malloch2009}. 

A communications protocol is a rigorous and formal description of a message format and rules for exchanging messages adhering to the protocol. The term protocol usually refers to the formal description of the protocol and is distinct from the implementation of the protocol. A protocol without an implementation is metaphorically like a constitution without a government. If two devices adhere to the same protocol then they are able to reliably communicate with each other because they can, in a sense, speak the same language. The mapping protocol is itself defined in terms of the Open Sound Control (OSC) protocol. This is a common practice for constructing increasingly specific protocols. The Internet is essentially a particular stack of protocols that are layered on top of each other and specify how devices and routers are to be addressed, open connections with each other, and interpret messages.

To say that one protocol is layered on top of another is to say that the higher level protocol is defined in terms of the structure and conventions of the lower level protocol. In the case of Libmapper, the protocol layer that the OSC protocol is implemented on top of is called the User Datagram Protocol (UDP). The other common protocol that can be used at this conceptual layer of the Internet protocol stack is the Transmission Control Protocol (TCP). The well-known Hypertext Transfer Protocol (HTTP) is layered on top of TCP. A web browser implements HTTP so that it can fetch and display web pages from remote systems that speak the same protocol. As a protocol, OSC in not defined in terms of a particular Transport protocol (the name given to the layer in the protocol that UDP and TCP occupy - HTTP occupies the layer on top of the Transport layer, known as the Application layer). This is different from HTTP, which is specifically a TCP-based protocol.

Libmapper then allows one to configure a mapping relationship between any collection of devices (control devices, sound generation devices, and anything else) by embedding Libmapper into the code for each device and creating an instance of Libmapper upon startup of each device, thus allowing any system on a local network to easily implement the mapping protocol without in depth knowledge of the protocol.

The specification of the OSC protocol can be found on the web \cite{osc2011}. In the opinion of the creators of the protocol \cite{osc2009}, OSC is a "weak" form of protocol because it does not define certain necessary aspects of any form of data transmission. It does not define the patterns of command and response, error handling or receipt acknowledgement - only the format of the messages that are sent. These aspects are left as decisions to be made by particular implementations of the protocol or a protocol stacked on top of OSC.

The complete definition of the OSC protocol is not important for the scope of this thesis, however the following aspects of the protocol should give a sufficient sense of what is defined:

\begin{description}
\item a unit of transmission is a \emph{OSC Packet}, which consists of a continguous block of binary data that is the contents of the packet and a number reflecting the number of bytes of data in the contents
\item the size of an \emph{OSC Packet} is always a multiple of 4 bytes
\item the contents of an \emph{OSC Packet} is an \emph{OSC Message}
\item an \emph{OSC Message} consists of an \emph{OSC Address Pattern} followed by an \emph{OSC Type Tag String} followed by zero or more \emph{OSC Arguments}
\item an \emph{OSC Address Pattern} is a URL-like string strictly beginning with the '/' character
\item an \emph{OSC Type Tag String} is a string strictly beginning with the ',' character followed a number of characters corresponding to the number of \emph{OSC Arguments} each of which specifies the OSC defined data type of each argument
\end{description}

\subsection{Mapper Protocol}

The semantics for the common set of OSC messages that the \emph{Mapper Protocol} (the name that will be used to describe the specification implemented by Libmapper in this paper) is intended to be as useful as possible for configuring mappings between output and input signals without assuming too much about what form the control or sound generation will take. In the Mapper Protocol, there is a concept of \emph{source device} and \emph{destination device}. In a DMI, the control mechanism will often be announced as a source device and the sound generation mechanism will be announced as a destination device. The set of messages proposed by the Mapper Protocol allows each device participating in the composite distributed system to \emph{announce} its presence, \emph{discover}, \emph{describe} itself, \emph{link} to another device, \emph{connect} one of its datastreams (also referred to as signals) to a datastream on a device it is linked to, and \emph{condition or transform} a signal with a basic mathematical function \cite{Malloch2009}. These last two pieces of functionality are the two possible levels of mapping that can be present in the prototypical mapping system referenced in section \ref{sec:Importance of Mapping}.

\subsubsection{Announcement}

Each device in the distributed system announces itself with a port and unique name that it can be addressed with that it would like to use. In the case that other devices in the local network have already reserved a particular port or unique name, Libmapper uses a collision detection algorithm to systematically try other ports and/or modify the requested name with an appended unique ordinal number. This allows a set of identical control or sound generation devices to each be programmed to request the same collection of OSC names for their signals without worrying about the identical devices being addressed by the same names.

\subsubsection{Discovery}

After a device appears on the network, announces itself, and receives a unique name and port number on the local network through the port and name allocation algorithm, it sends the following message to request that all Libmapper-compatible devices describe themselves:

\begin{quote}
\verb#/who#
\end{quote}

Other devices are obliged by the protocol to respond with (for example):

\begin{quote}
\verb#/registered /tstick.1 @inputs 1 @outputs 52 @class /tstick# 

\verb#  @IP 192.168.0.3 @port 8001#

\verb#/registered /granul8.1 @inputs 80 @outputs 0 @class /granul8#

\verb#  @IP 192.168.0.4 @port 8000#
\end{quote}

In this way any device or monitor (in the case of a graphical user interface) can request the names, IP addresses, and UDP ports of any devices on the network. 

\subsubsection{Device Linking}

To create a direct network link between two devices, the Mapper Protocol specifies the following message:

\begin{quote}
\verb#/link /tstick.1 /granul8.1#
\end{quote}

Libmapper creates a router in software, which can be thought of as a special virtual device distinct from source devices and destination devices, for each pair of linked devices. In Mapper Protocol terminology, a link is created between a source device and a destination device. Every router is tasked with the signal conditioning and message transformation for any pair of connected signals involving the two devices in the source/destination link pair that the specific router is created by Libmapper to manage. This Libmapper design decision allows each device acting as a source in a link to handle the traffic and processing of the links that it is involved in, thus distributing the work load over all source devices as opposed to assuming a single central router handling all traffic between all linked devices in the system.

When the Libmapper instance in the source device sucessfully creates a router for the link and completes initialization of the connection after receiving a \verb#/link# message, it responds with:

\begin{quote}
\verb#/linked /tstick.1 /granul8.1#
\end{quote}

Two devices are unlinked by sending the following message:

\begin{quote}
\verb#/unlink /tstick.1 /granul8.1#
\end{quote}

The source device then destroys the router handling this link after receiving a \verb#/unlink# message, and responds with:

\begin{quote}
\verb#/unlinked /tstick.1 /granul8.1#
\end{quote}

\subsubsection{Datastream Connection}

Once a router has been created between a source and destination by requesting a link, specific output signals on the source device can be connected to specific input signals on the destination device. For example, once a T-Stick (\verb#/tstick.1#) has been linked to the Granul8 granular synthesizer (\verb#/granul8.1#), the individual signals output by the touch sensors, accelerometers, etc. in the T-Stick can be routed to specific inputs in the granular synthesizer. Without specifying connections between output signals and input signals then gestures performed with the control source device will have no affect on the audio samples produced by the sound generation destination device. 

In the Mapper Protocol, this is accomplished by sending a message similar to these:

\begin{quote}
\verb#/connect /tstick.1/raw/pressure /granul8.1/gain#

\verb#/connect /tstick.1/raw/pressure /granul8.1/gain#

\verb#  @scaling expression @expression x*10 @clipping minimum 0#
\end{quote}

The router managing the link between \verb#/tstick.1# and \verb#/granul8.1# records all the relevant information communicated by the \verb#/connect# message, sets up the address translation and transformation required to connect the output signal to its associated input signal, and responds with:

\begin{quote}
\verb#/connected /tstick.1/raw/pressure /granul8.1/gain#

\verb#/properties /tstick.1/raw/pressure /granul8.1/gain#

\verb#  @scaling expression#

\verb#/properties /tstick.1/raw/pressure /granul8.1/gain#

\verb#  @expression x*10#

\verb#/properties /tstick.1/raw/pressure /granul8.1/gain#

\verb#  @clipping minimum 0#
\end{quote}

A signal pair is disconnected by sending:

\begin{quote}
\verb#/disconnect /tstick.1/raw/pressure /granul8.1/gain#
\end{quote}

The router responds with:

\begin{quote}
\verb#/disconnected /tstick.1/raw/pressure /granul8.1/gain#
\end{quote}

The properties of a connection can be queried with:

\begin{quote}
\verb#/connection/properties/get /tstick.1/raw/pressure /granul8.1/gain#
\end{quote}

The syntax of the message for modifying properties of a connection is similar to creating a connection with \verb#/connect# replaced by \verb#/connection/modify#.

These special OSC messages as defined by their message headers and syntax comprise what is referred to in this paper as the \emph{Mapper Protocol}. By sending messages from this set of messages recognized by Libmapper it is possible to configure mappings between an arbritrarily large collection of devices on a local network with relatively little work.

\section{Sensor Networks}

A \emph{sensor network} is a term used in electrical engineering literature to refer to a system that employs a (typically) large number of distributed sensors that are connected to each other through autonomous, intermediary nodes capable of processing and communicating the data obtained from the particular sensors it is attached to and relaying that information to other nodes on the network. If the nodes communicate with each other through some form of wireless radio then the network can be referred to as a \emph{wireless sensor network}. The key functional difference being that in the wireless scenario, the nodes are not tethered to a particular physical location and are free to move provided they remain within radio range. Libmapper can, in principle, with few to no modifications operate in either scenario and the phrase \emph{sensor network} will be used to encapsulate both.

The concept of a sensor network is useful because it strikes a good balance of specificity and abstraction when describing a particular network topology and implementation. Sensor networks, particularly wireless sensor networks, are an active field of research in a diverse array of research fields including defense, environmental monitoring, energy/heat systems, and structural engineering. In these and related situations, the function of the sensor network is to sense, measure, and gather information from the environment and/or actors in the environment and use the capability of distributed processing nodes to make decisions about where the information should be routed \cite{survey2008}. The majority of research applications do not focus on the use of these datastreams to output sound, visuals, or control actuators in realtime, however the definition of a sensor network does not preclude the nodes in the network from being connected to actuators as well as sensors.

\section{Digital Musical Instruments as Localized Sensor Networks}

A DMI can be understood as a particular form of sensor network that is particularly localized and stable. Particularly, if all of the sensors, processing nodes (perhaps each running an instance of Libmapper), and sound generation mechanism are all housed in a single physical object. Typically, no academic would refer to such a DMI as a sensor network, however it would technically fall under the working definition used in this paper. A sensor network may include functionality that compensates for unreliable sensors and devices that enter and leave the network relatively often, especially for a large, distributed network, which is not necessary in a DMI with comparatively few and robust sensors. The larger the network, the more effort must be made to handle sensor and device failure. Hardware failure and reliability is a topic is outside of the scope of this thesis, but is an important practical distinguish between systems that are typically termed \emph{DMIs} and systems termed \emph{sensor networks}.

Framing a DMI as a sensor network affords one the ability to imagine what form a more \emph{sensor network-like} DMI might take. Additionally, because the task of mapping is fundamentally seen as a task of network topology configuration in this thesis, it is a better term to use when focusing on the mapping layer of a musical system. A single DMI might consist of many spatially-distributed control devices that each output signals connected to the inputs of a sound generation device. Or an orchestra of DMIs might be understood to be one musical sensor network that together produce one collection of sound. The argument is that understanding what is being built in terms of a sensor network, as opposed to individual DMIs that may or may not be played together, better reflects the pecularities of what a DMI fundamentally is and why the task of building a DMI is so different from building an acoustic musical instrument.

Additionally, perhaps such a reframing might better work \emph{with} the pecularities of a DMI and lessen the extent to which design teams work against the inate characteristics of computers attempting to make DMIs in the image of acoustic musical instruments. After all, as previously discussed, this fundemental difference between a DMI and acoustic musical instrument is what introduces the task of mapping as unique to DMIs in the first place. For convenience, the term \emph{musical sensor network} (MSN) will be used in the remainder of this paper to reference this particular framing of the DMI concept in contexts that might be helped by this alternative concept. However, the two terms can usually be used interchangably without being inaccurate. 

Maxmapper vs. vizmapper and DMI vs. sensor network

\section{Usability and Graphical User Interfaces}

Imagine a MSN comprised of more than ten Libmapper-enabled devices, some of them are control devices and some are sound generation devices. Each device might output, or receive as inputs, around ten signals. So to create a mapping layer over this distributed network of devices, one might send hundreds of Mapper Protocol messages to link, unlink, connect, disconnect, etc. various pairs of devices and signals as one attempts to determine through trial and error the specific set of connections and functional transformations that best fit the particular sensors, devices, context, and design intentions of the people constructing the system. 

It may be that it is unclear what control device ought to be linked to what sound generation device, much less what signals within the control device to connect to what particular signals within the sound generation device. When experimenting with a mapping layer for a specific MSN it is necessary to view the current network mapping, modify the current mapping, test the new mapping, evaluate its effectiveness and then repeat the process until the mapping is satisfactory. 

Any tools and interfaces that speed up this iterative process and reduce the amount of Libmapper-specific/Mapper Protocol knowledge required to efficiently explore the possibilities of a particular MSN improves the usability of Libmapper and allows for a wider diversity of skill sets and roles in a broader performance project involving instrumentalists, composers, directors, and engineers to participate in the mapping task that is so crucial to the character of a musical system. 

Advances in user interface design are historically a need driven by the specifics of the nature of the system that the user requires an interface to manipulate. Similarly, data visualization design is driven by the need of a user to develop understanding of a system from a complex set of data about a system so that they can correctly manipulate the system to match their intentions. This chapter has been a foundation for understanding the interface and visualization design decisions that are made in chapter 4 as coupled to the requirements and need for a mapping system. The next chapter presents a foundation for understanding the decisions in chapter 4 as coupled to what is known about what type of visualizations make a system easy to understand and what type of interfaces are effective for manipulating systems.
